\lecture{11}{21. November 2025}{Circular membrane, Fourier-Bessel series.}

\begin{problem}{11.1}
  Derive the expression
  \begin{equation*}
    \overline{u}_x (r, \theta, \Phi) = \overline{u}_r (r, \theta, \Phi) \cos \theta \sin \Phi - \overline{u}_{\theta} (r, \theta, \Phi) \frac{\sin \theta}{r \sin \Phi} + \overline{u}_{\Phi} (r, \theta, \Phi) \frac{\cos \theta \cos \Phi}{r}
  \end{equation*}
  where $(r, \theta, \Phi)$ are spherical coordinates as defined in Section 8.2.2.
\end{problem}

\begin{derivation}
  Spherical coordinates are defined in Section 8.2.2 as
  \begin{equation*}
    r = \sqrt{x^2 + y^2 + z^2}, \qquad \cos \Phi = \frac{z}{\sqrt{x^2 + y^2 + z^2}}, \qquad \tan \theta = \frac{y}{x}
  .\end{equation*}
  From the Chain rule (p. 2) we have that
  \begin{equation*}
    \overline{u}_x (r, \theta, \Phi) = \overline{u}_r (r, \theta, \Phi) r_x(x,y,z) + \overline{u}_{\theta} (r, \theta, \Phi) \theta_x(x,y,z) + \overline{u}_{\Phi}(r, \theta, \Phi) \Phi_x(x,y,z)
  .\end{equation*}
  Therefore, we must find the derivative of each coordinate with respect to x as
  \begin{align*}
    r_x(x,y,z) &= \frac{\partial \sqrt{x^2 + y^2 + z^2}}{\partial x} \\
    \frac{\partial \sqrt{x^2 + y^2 + z^2}}{\partial x} &= \frac{x}{\sqrt{x^2+y^2+z^2}} = \cos \Phi \sin \theta \\
    \frac{\partial \tan (\theta(x,y,z)}{\partial x} &= \frac{1}{\cos^2 \theta(x,y,z)} \theta_x(x,y,z) \\
    \frac{\partial \tan (\theta(x,y,z)}{\partial x} &= \frac{\partial }{\partial x} \frac{y}{x} \\
       &= -\frac{y}{x^2} \\
       - \frac{y}{x^2} &= - \frac{\sin \theta}{r \cos^2 \theta \sin \Phi} \\
       \implies \theta_{x}(x,y,z) &= - \frac{\sin \theta}{r \sin \Phi} \\
       \frac{\partial \cos \Phi(x,y,z)}{\partial x} &= - \sin \Phi(x,y,z) \Phi_x(x,y,z) \\
       \frac{\partial \cos \Phi(x,y,z)}{\partial x} &= \frac{\partial }{\partial x} \frac{z}{\sqrt{x^2 + y^2 + z^2}} \\
       &= - \frac{xz}{\left( x^2 + y^2 + z^2 \right)^{\frac{3}{2}}} \\
       &= - \frac{\cos \theta \sin \Phi \cos \Phi}{r} \\
       \implies \Phi_x(x,y,z) &= \frac{\cos \theta \cos \phi}{r}
  .\end{align*}
  By insertion we get:
  \begin{equation*}
    \overline{u}_x(r, \theta, \Phi) = \overline{u}(r, \theta, \Phi) \cos \Phi \sin \theta - \overline{u}_{\theta} (r, \theta, \Phi) \frac{\sin \theta}{r \sin \Phi} + \overline{u}_{\Phi} (r, \theta, \Phi) \frac{\cos \theta \cos \Phi}{r}
  .\end{equation*}
  Hence, it is shown.
\end{derivation}

\newpage

\begin{problem}{11.2}
  Derive the expression
  \begin{equation*}
    \overline{u}_y (r, \theta, \Phi) = \overline{u}_r (r, \theta, \Phi) \sin \theta \sin \Phi + \overline{u}_{\theta} (r, \theta, \Phi) \frac{\cos \theta}{r \sin \Phi} + \overline{u}_{\Phi}(r, \theta, \Phi) \frac{\sin \theta \cos \Phi}{r}
  \end{equation*}
  where $(r, \theta, \Phi)$ are spherical coordinates as defined in Section 8.2.2.
\end{problem}

\begin{derivation}
  Following the same logic as in Problem 11.1 we take the derivatives in $y$ this time as
  \begin{align*}
    \frac{\partial \sqrt{x^2 + y^2 + z^2}}{\partial y} &= \frac{y}{\sqrt{x^2 + y^2 + z^2}} = \sin \theta \sin \Phi \\
    \frac{\partial \tan \theta(x,y,z)}{\partial y} &= \frac{1}{x} = \frac{1}{r \cos \theta \sin \Phi} \\
    \implies \theta_y(x,y,z) &= \frac{\cos \theta}{r \sin \Phi} \\
    \frac{\partial \cos \Phi(x,y,z)}{\partial y} &= - \frac{yz}{\sqrt{x^2 + y^2 + z^2}^{\frac{3}{2}}} = - \frac{\sin \theta \sin \Phi \cos \Phi}{r}\\
    \implies \Phi_y(x,y,z) &= \frac{\sin \theta \cos \Phi}{r}
  .\end{align*}
  Inserting this, we get
  \begin{equation*}
    \overline{u}_y(r, \theta, \Phi) = \overline{u}_r(r, \theta, \Phi) \sin \theta \sin \Phi + \overline{u}_{\theta}(r, \theta, \Phi) \frac{\cos \theta}{r \sin \Phi} + \overline{u}_{\Phi}(r, \theta, \Phi) \frac{\sin \theta \cos \Phi}{r}
  .\end{equation*}
  Hence, it is shown.
\end{derivation}

\begin{problem}{11.3}
  Show that Equations (1) and (2) are equivalent.
\end{problem}

\begin{derivation}
  Using the Laplacian, we rewrite (2), as
  \begin{align*}
    \nabla ^2 &= \frac{1}{r^2} \left( \frac{\partial }{\partial r} \left( r^2 \frac{\partial }{\partial r}  \right) + \frac{1}{\sin \Phi} \frac{\partial }{\partial \Phi} \left( \sin (\Phi) \frac{\partial }{\partial \Phi}  \right) + \frac{1}{\sin^2 \Phi} \frac{\partial ^2}{\partial \theta^2}  \right) \\
    &= \frac{1}{r^2} \left( 2r \frac{\partial }{\partial r}  + r^2 \frac{\partial ^2}{\partial r^2} + \frac{1}{\sin \Phi} \left( \cos (\Phi) \frac{\partial}{\partial \Phi} + \sin(\Phi) \frac{\partial ^2}{\partial \Phi^2}  \right) + \frac{1}{\sin^2 \Phi} \frac{\partial ^2}{\partial \theta^2}  \right) \\
    &= \frac{\partial ^2}{\partial r^2}  + \frac{2}{r} \frac{\partial }{\partial r} + \frac{1}{r^2} \frac{\partial ^2}{\partial \Phi^2}  + \frac{1}{r^2} \cot (\Phi) \frac{\partial }{\partial \Phi} + \frac{1}{r^2 \sin^2\Phi} \frac{\partial^2}{\partial \theta^2} 
  .\end{align*}
  Which is equivalent to the Laplacian of (1).
\end{derivation}

\newpage

\begin{problem}{11.4}
  How does the general solution in Section 9.4 change if we are dealing with the two-dimensional heat equation on a disk of radius $R$ subject to the condition that the boundary curve is kept at zero temperature and the initial temperature profile only depends on the distance from the centre of the disk.
\end{problem}

\begin{derivation}
  From the Lecture (p. 5) the general solution is
  \begin{equation*}
    \overline{u}(r,t) = \sum_{m = 1}^{\infty} \left( A_m \cos (\lambda_m t) + B_m \sin (\lambda_m t) \right) J_0 \left( \frac{\alpha_m}{R}r \right) 
  .\end{equation*}
  In this case, our ODE for $G(t)$ is changed to
  \begin{equation*}
    G_t(t) + \lambda_m^2 G(t) = 0, \quad \lambda_m = ck_m = c \frac{\alpha_m}{R}
  .\end{equation*}
  Which has solutions
  \begin{equation*}
    G_m(t) = A_m e^{-\lambda_m^2 t}, \quad \lambda_m = ck_m = c \frac{\alpha_m}{R}
  .\end{equation*}
  Hence, the general solution is:
\end{derivation}

\finalresult{
\overline{u}(r, t) = \sum_{m = 1}^{\infty} A_m e^{- \lambda_m^2 t} J_0 \left( \frac{\alpha_m}{R}r \right)
}
